\abstract{
The Internet of Things has instrumented billions of homes, buildings, farms, and factories
with smart devices, yet today's platforms fail to meet the needs of the people who rely on
them. Control logic is concentrated in vendor clouds and hubs that add latency, cost, and
single points of failure; fire-and-forget RPCs leave developers and occupants without
feedback on whether actions take effect; and users navigate a rich control spectrum with no
structured support. This thesis argues that smart-space control should be designed around the
needs of the people who inhabit the space, across three planes of the smart-space
stack---the \emph{action}, \emph{control}, and \emph{management} planes.

Our work straddles distributed systems and human-computer interaction. At the
\emph{action plane}, RASC (Request-Ack-Start-Complete) replaces fire-and-forget RPC with an
abstraction that exposes the full action lifecycle atop existing platforms, so control logic
and occupants can observe and program what their devices do; proofs and experiments confirm
fast, low-overhead lifecycle detection and large scheduling gains over state-of-the-art
baselines. At the \emph{control plane}, CoMesh is the first fully decentralized, hub-less
Sense-Trigger-Actuate control plane for large commercial edge IoT meshes; using group-based
coordination with zero-message exchange and locality-sensitive hashing, it keeps control
local, scalable, and fault-tolerant---removing the hub bottlenecks and cloud dependency that
cost occupants latency, money, and uptime. At the \emph{management plane}, Control in Context,
a study of smart home owners (survey and interviews), shows that control is a fluid,
situational process shaped by reliability and cognitive load rather than ecosystem
boundaries, grounding interface design in how people actually behave.

Though independent, the three contributions share one vision: each makes a plane of
smart-space control resilient, observable, and built around real user needs.
}